% Standalone problem document, generated from the adjacent README.md.
% Compile with XeLaTeX. Mathematical content is maintained in Markdown.
\documentclass[11pt,a4paper]{article}
\usepackage[fontsize=10.5pt]{scrextend}
\usepackage[margin=22mm,headheight=15pt,headsep=9mm,footskip=11mm]{geometry}
\usepackage{fontspec}
\setmainfont{texgyrepagella}[Extension=.otf,UprightFont=*-regular,BoldFont=*-bold,ItalicFont=*-italic,BoldItalicFont=*-bolditalic]
\setsansfont{texgyreheros}[Extension=.otf,UprightFont=*-regular,BoldFont=*-bold,ItalicFont=*-italic,BoldItalicFont=*-bolditalic]
\setmonofont{texgyrecursor}[Extension=.otf,UprightFont=*-regular,BoldFont=*-bold,ItalicFont=*-italic,BoldItalicFont=*-bolditalic,Scale=MatchLowercase]
\usepackage{amsmath,amssymb}
\usepackage{unicode-math}
\setmathfont{texgyrepagella-math.otf}
\usepackage{xcolor,microtype,enumitem,needspace,fancyhdr}
\usepackage{bookmark}
\definecolor{ink}{HTML}{17344B}
\definecolor{accent}{HTML}{197D80}
\definecolor{muted}{HTML}{596773}
\hypersetup{colorlinks=true,linkcolor=accent,urlcolor=accent,pdftitle={MF-23 - Complete
Crouzeix
conjecture},pdfauthor={Open Problems in Numerical Linear Algebra}}
\urlstyle{same}
\setlength{\parindent}{0pt}
\setlength{\parskip}{5pt plus 1pt minus 1pt}
\linespread{1.04}
\setlength{\emergencystretch}{3em}
\widowpenalty=10000
\clubpenalty=10000
\displaywidowpenalty=10000
\allowdisplaybreaks[1]
\setlist{nosep,leftmargin=1.5em}
\providecommand{\tightlist}{\setlength{\itemsep}{0pt}\setlength{\parskip}{0pt}}
\providecommand{\pandocbounded}[1]{#1}
\pagestyle{fancy}
\fancyhf{}
\fancyhead[L]{\sffamily\footnotesize\color{muted}OPEN PROBLEMS IN NUMERICAL LINEAR ALGEBRA}
\fancyhead[R]{\sffamily\bfseries\footnotesize\color{accent}MF-23}
\fancyfoot[L]{\parbox[b]{0.9\textwidth}{\sffamily\fontsize{8}{10}\selectfont\color{muted}Editorial ratings; a bounded search cannot certify that a problem remains unsolved.\\Literature check: 2026-09-11}}
\fancyfoot[R]{\sffamily\footnotesize\color{muted}\thepage}
\renewcommand{\headrulewidth}{0.3pt}
\renewcommand{\headrule}{\hbox to\headwidth{\color{accent}\leaders\hrule height \headrulewidth\hfill}}
\makeatletter
\renewcommand{\section}{\Needspace{5\baselineskip}\@startsection{section}{1}{0pt}{12pt plus 2pt minus 2pt}{4pt}{\sffamily\bfseries\large\color{ink}}}
\renewcommand{\subsection}{\Needspace{4\baselineskip}\@startsection{subsection}{2}{0pt}{9pt plus 1pt minus 1pt}{3pt}{\sffamily\bfseries\normalsize\color{ink}}}
\makeatother
\setcounter{secnumdepth}{0}
\begin{document}
{\sffamily\small\color{muted}Matrix functions and stability\par}
\vspace{3pt}
{\raggedright\hyphenpenalty=10000\exhyphenpenalty=10000\sffamily\bfseries\fontsize{21}{24}\selectfont\color{ink}Complete
Crouzeix conjecture\par}
\vspace{7pt}
{\sffamily\small\textbf{Difficulty:} extreme\quad\textbf{Importance:} broadly
interesting\par}
{\sffamily\small\bfseries\color{accent}Status: Partially resolved\par}
\vspace{4pt}
{\color{accent}\hrule height 0.8pt}
\vspace{6pt}
\textbf{Topic:} Matrix-valued polynomial functional calculus

\textbf{Rating rationale:} The uniform bound must survive arbitrary
matrix amplification, beyond the scalar functional calculus; it would
give sharp bounds for block matrix functions and connect numerical-range
estimates to dilation and similarity theory.

\subsection{Context and notation}\label{context-and-notation}

All matrices are complex, and all norms are Euclidean operator norms.
For \(A\in\mathbb C^{n\times n}\) define

\[
W(A)=\{x^*Ax:x\in\mathbb C^n,\ x^*x=1\}.
\]

If \(F(z)=[f_{ij}(z)]_{i,j=1}^m\) has polynomial entries, let
\(F(A)=[f_{ij}(A)]_{i,j=1}^m\), an \(mn\times mn\) block matrix.

\subsection{Problem statement}\label{problem-statement}

Prove or disprove that, for every \(n,m\geq1\), every
\(A\in\mathbb C^{n\times n}\) and every matrix polynomial
\(F\in\mathbb C^{m\times m}[z]\),

\[
\|F(A)\|_2\leq 2\max_{z\in W(A)}\|F(z)\|_2.
\]

This is the complete, rather than scalar, numerical-range spectral-set
question. It controls simultaneous polynomial functions and hence block
approximation errors with one constant independent of matrix and block
sizes.

\subsection{References}\label{references}

\begin{itemize}
\tightlist
\item
  Michel Crouzeix,
  \href{https://doi.org/10.1016/j.jfa.2006.10.013}{\emph{Numerical range
  and functional calculus in Hilbert space}}, \emph{Journal of
  Functional Analysis} \textbf{244}(2), 668--690 (2007). Original source
  for the complete conjecture, as explicitly attributed in the next
  source.
\item
  Per Åhag, Rafał Czyż and Jani Virtanen,
  \href{https://arxiv.org/html/2608.27346v3}{\emph{Square Functions,
  Complete Crouzeix Conjecture in Dimension Three, and the
  Clouâtre--Ostermann--Ransford conjecture}, arXiv:2608.27346v3},
  \textbf{9 September 2026}, Introduction, equation \textbf{(1.2)} and
  \textbf{Theorem 1.1}. This current version states the exact target and
  explicitly distinguishes it from the scalar proofs. It claims the
  complete bound for \(n\leq3\). Its counterexample reduction leaves
  only \(2\times2\) matrix polynomials to check when \(n=4\).
\item
  Michel Crouzeix and César Palencia,
  \href{https://doi.org/10.1137/17M1116672}{\emph{The numerical range is
  a \((1+\sqrt2)\)-spectral set}}, \emph{SIAM Journal on Matrix Analysis
  and Applications} \textbf{38}(2), 649--655 (2017). The complete
  universal upper bound is \(1+\sqrt2\).
\end{itemize}

\subsection{Scope and status check}\label{scope-and-status-check}

The target is source-stated, not an editorial extension of the scalar
conjecture. The August 2026 scalar proof claims by Jin and by
Lorist--Schwenninger do not settle it: the September 9 primary
manuscript explicitly addresses this distinction. The \(n=3\) result is
a recent preprint claim; the full displayed target remains unresolved
even if that claim is accepted. Searches through the check date for
complete Crouzeix proofs, matrix-valued Crouzeix inequalities and later
versions found no full solution. The version was inspected as both HTML
and PDF; the inequality is on PDF page 2, immediately before Theorem
1.1.
\end{document}
